\chapter{Conclusion and Future Work}
\label{chap:conclusion}

\section{Summary of Completed Work}
Urban cycling infrastructure planning has historically suffered from an empirical blind spot: municipal traffic authorities rely almost exclusively on retrospective police accident records, which register only catastrophic collisions resulting in bodily injury or significant property damage. The vast majority of hazardous vehicular interactions---such as aggressive close-passing maneuvers, high-speed aerodynamic buffeting, and near-miss emergency swerves---go completely unrecorded. Furthermore, existing naturalistic cycling initiatives either lack objective lateral distance transducers or evaluate proximity through crude, binary thresholding that ignores vehicle speed differentials, vehicle mass disparities, and cornering tilt dynamics.

This project presented \textbf{CycleSafe}, a holistic, dual-domain research and engineering platform combining physical IoT telemetry and high-resolution microscopic traffic simulation to quantify, classify, and visualize urban cyclist overtaking safety. The completed work spans four core pillars:
\begin{enumerate}
    \item \textbf{Instrumented Hardware Telemetry Design:} An open-source embedded sensor node was designed based on the dual-core ESP32 microcontroller, dual waterproof ultrasonic transducers (JSN-SR04T), u-blox NEO-6M satellite positioning, and an MPU-6050 six-axis inertial measurement unit. We formulated an analytical roll-angle tilt compensation model ($\phi = \arctan(a_y / \sqrt{a_x^2 + a_z^2})$) that mathematically projects slant Time-of-Flight distances onto the true horizontal plane ($d_{\text{true}} = d_{\text{measured}} \cos\phi$), correcting for cyclist lean during cornering and pedaling sway. A dedicated 3.3\,V AMS1117 regulator rail was integrated to isolate high-current SD SPI write spikes, preventing flash filesystem corruption.
    \item \textbf{Microscopic Simulation Digital Twin:} Using Eclipse SUMO and TraCI, we established a calibrated, multi-modal simulation of a 9.53\,km arterial thoroughfare connecting Berlin's Ernst-Reuter-Platz to Oberbaumbr{\"u}cke. Driven by empirical 12-hour vehicular census data (13,531 motor vehicles and 1,397 bicycles), the simulation executed 7,201 discrete sublane time steps per 60-minute run across 1,244 concurrent vehicles, logging sub-second spatial envelopes and kinematic velocities across passenger cars, heavy freight trucks, motorcycles, and commuter bicycles.
    \item \textbf{Machine Learning Conflict Severity Engine:} Moving beyond simplistic distance cutoffs, an Extreme Gradient Boosting (XGBoost) model was trained on the TU Berlin SimRA naturalistic cycling dataset. By synthesizing lateral clearance, relative overtaking velocity, deceleration shock, and counterpart vehicle kinetic mass into a multi-class objective function, the model achieves a 90.6\% F1-score in classifying encounters into \textit{Safe Pass}, \textit{Close Pass (Medium)}, and \textit{Severe Incident}, outperforming static German StVO \S\,5(4) distance rules by 16.2 percentage points.
    \item \textbf{Full-Stack Geospatial Web Interface:} A containerized web platform orchestrated via Docker Compose was engineered using Go, PostgreSQL, Redis, Leaflet, and Chart.js. The interface supports interactive multi-bicycle trajectory switching across 15 tracked bicycles, spatial conflict clustering, longitudinal kinematic acceleration plotting, and RFC 4180 compliant telemetry CSV export.
\end{enumerate}

\section{Achievement of Objectives}
The five concrete objectives formulated in Section \ref{sec:aim_objectives} have been successfully fulfilled, as evidenced in Table \ref{tab:objective_fulfillment} and the referenced report sections.

\begin{table}[h!]
	\centering
	\caption{Systematic mapping and verification of project objectives.}
	\label{tab:objective_fulfillment}
	\setlength{\tabcolsep}{5pt}
	\renewcommand{\arraystretch}{1.25}
	\begin{tabularx}{\textwidth}{|>{\raggedright\arraybackslash\bfseries}p{2.8cm}|>{\raggedright\arraybackslash}X|>{\raggedright\arraybackslash}p{4.2cm}|>{\raggedright\arraybackslash}p{2.6cm}|}
		\hline
		\textbf{Objective} & \textbf{Key Deliverables \& Specs} & \textbf{Quantitative Evidence} & \textbf{Status} \\
		\hline
		\hline
		\textbf{Obj 1:} Embedded Hardware & ESP32, dual JSN-SR04T, NEO-6M, MPU-6050, AMS1117 LDO, tilt math. & Rail stability $\pm 18\,\text{mV}$; roll tilt error $<1.0\%$ at $30^\circ$. & \textbf{Achieved} (Sec. \ref{sec:tilt_comp}, \ref{sec:hardware_bom}, \ref{sec:hardware_benchmarks}) \\
		\hline
		\textbf{Obj 2:} Mechanical Packaging & Weatherproof PETG enclosure and handlebar mount for FDM print. & 14-item BOM: INR~3,494 to Dhanbad; $T_g \approx 85^\circ\text{C}$ heat stability. & \textbf{Achieved} (Sec. \ref{sec:enclosure_specs}, \ref{sec:hardware_bom}, App. \ref{app:bom}) \\
		\hline
		\textbf{Obj 3:} Simulation Twin & 10\,km dual-corridor in SUMO / TraCI ($\Delta y = 0.2\,\text{m}$); 1,244 vehicles. & 7,201 TraCI steps; 1,048 close proximity encounters evaluated. & \textbf{Achieved} (Sec. \ref{sec:simulation_setup}, \ref{sec:functional_validation}) \\
		\hline
		\textbf{Obj 4:} ML Conflict Classifier & Supervised XGBoost classifier on SimRA with speed \& vehicle mass. & Accuracy: 92.1\%; F1: 90.6\% (vs. 74.4\% for static 1.5\,m rule). & \textbf{Achieved} (Sec. \ref{sec:xgboost_formulation}, \ref{sec:ml_evaluation}) \\
		\hline
		\textbf{Obj 5:} Full-Stack Dashboard & Containerized Go, PostgreSQL, Redis, Leaflet maps, and CSV export. & Dynamic switching across 15 bicycles, heatmaps, RFC 4180 CSV export. & \textbf{Achieved} (Sec. \ref{sec:dev_environment}, \ref{sec:functional_validation}) \\
		\hline
	\end{tabularx}
\end{table}

\section{Main Contributions}
The primary scientific, methodological, and engineering contributions delivered by this research include:
\begin{enumerate}
    \item \textbf{Analytical Tilt Compensation for Naturalistic Bicycle Telemetry:} Conventional bicycle-mounted ultrasonic transducers suffer substantial geometric distortion when the rider leans into a turn or rocks laterally during vigorous uphill pedaling. By fusing three-axis gravitational vector components from the MPU-6050 accelerometer into an analytical trigonometric correction ($d_{\text{true}} = d_{\text{measured}} \cos\phi$), CycleSafe eliminates false close-pass triggers caused by road camber or cornering dynamics.
    \item \textbf{Empirically Grounded Multi-Modal Microscopic Digital Twin:} Unlike theoretical simulation models that treat bicycles as point masses or miniature automobiles, CycleSafe implements TraCI-driven lateral sublane dynamics with non-lane-centered maneuvering. The model incorporates distinct vehicle length, width, and acceleration profiles across four vehicle classes, generating realistic emergent overtaking behaviors calibrated against metropolitan census counts.
    \item \textbf{Vehicle Mass-Weighted Machine Learning Conflict Assessment:} While existing statutory passing distance laws enforce a uniform 1.5\,m rule regardless of vehicle size, CycleSafe establishes that kinetic risk scales quadratically with overtaking speed and linearly with counterpart vehicle mass ($E_k = \frac{1}{2} m v^2$). Our XGBoost classifier demonstrates superior discriminative ability, correctly identifying high-hazard heavy vehicle overtakes that fall just outside legal clearance thresholds while recognizing that close passes between peer bicycles on segregated paths are routine and safe.
    \item \textbf{Reproducible, Open-Source Full-Stack Architecture:} All simulation scripts, embedded firmware configurations, 3D CAD parameters, database schemas, and microservice container definitions are released under permissive open-source licenses, ensuring full reproducibility for academic institutions and civil transit authorities.
\end{enumerate}

\section{Technical Limitations}
To ensure academic rigor, several engineering constraints and technical limitations of the current implementation are acknowledged:
\begin{enumerate}
    \item \textbf{Acoustic Beam Dispersion at Steep Grazing Angles:} The JSN-SR04T ultrasonic transducer possesses a conical acoustic emission beam of approximately $45^\circ$. When overtaking vehicles possess aerodynamic or curved side contours (e.g., streamlined sports cars or angled truck under-guards), acoustic reflection angles may exceed the transducer's acceptance envelope, resulting in sporadic pulse loss or echo attenuation.
    \item \textbf{Planar Corridor Approximation in Microscopic Simulation:} The SUMO simulation corridor is modeled with flat topographical gradients ($0\%$ incline). While accurate for Berlin's predominantly level road network, applying the model to undulating terrain (such as Indian plateau cities) would require incorporating gravitational gradient effects on cyclist pedaling speed and heavy vehicle braking distances.
    \item \textbf{Urban Canyon Satellite Multipath Effects:} The u-blox NEO-6M GPS receiver utilizes standard L1 C/A civil satellite signals with a circular error probable (CEP) of 2.5\,m in open-sky conditions. In dense high-rise urban street canyons, satellite multipath reflections can induce temporary positional drift, necessitating map-matching algorithms for trajectory alignment.
    \item \textbf{Uncalibrated Mixed-Traffic Heterogeneity for Developing Regions:} The microscopic simulation parameters are calibrated against European urban driving standards with high lane discipline. In non-lane-based traffic environments typical of Indian cities (e.g., high densities of three-wheelers, auto-rickshaws, and non-motorized cycle rickshaws), lateral maneuvering behaviors are significantly more chaotic and require localized driving behavior calibration.
\end{enumerate}

\section{Future Research Directions}
Building upon the verified foundation of the CycleSafe platform, several promising directions are outlined for subsequent phases of research and development:
\begin{enumerate}
    \item \textbf{Physical Field Trials on Dhanbad Road Network:} Following procurement and assembly of the verified INR~3,494 BOM, instrumented road trials will be conducted across selected arterial routes surrounding the IIT (ISM) Dhanbad campus (e.g., Sardar Patel Nagar Road, Hirapur Market corridor, and Grand Trunk Road / NH-19). Real-world sensor logs will provide comparative validation against the simulation predictions.
    \item \textbf{Integration of Edge Computer Vision for Counterpart Verification:} Replacing or augmenting ultrasonic sensors with low-power optical edge AI (e.g., an ESP32-CAM or Raspberry Pi Zero 2W running MobileNet-SSD) will enable real-time visual counterpart classification, allowing the system to extract vehicle license plates and classify vehicle types directly from edge video streams.
    \item \textbf{Automated Full-City OpenStreetMap Network Extraction:} Extending the simulation framework from the current 10\,km dual-corridor to an automated OSM graph ingestion pipeline that dynamically compiles entire metropolitan networks, including traffic signals, roundabouts, and multi-lane intersections.
    \item \textbf{V2X Connected Vehicle Communication:} Integrating low-latency wireless communication (such as Cellular-V2X or DSRC 802.11p) to broadcast acoustic proximity alerts directly to approaching vehicles, enabling automated collision avoidance warnings.
\end{enumerate}
